% =====================================================
% CHAPITRE 6 : ÉVALUATION ET INTERPRÉTABILITÉ (CRISP-DM Phase 5)
% =====================================================
\chapter{Évaluation et interprétabilité}

% ---------------------------------------------------
\section*{Introduction}
\addcontentsline{toc}{section}{Introduction}
% ---------------------------------------------------

Avoir un modèle qui fonctionne sur les données d'entraînement ne suffit pas. L'évaluation
sert à répondre à deux questions distinctes : le modèle détecte-t-il réellement les fraudes
(évaluation quantitative sur métriques adaptées au déséquilibre) ? Et ses décisions
sont-elles compréhensibles et défendables pour un auditeur (évaluation qualitative via
SHAP, LIME et Ollama) ? Ces deux dimensions sont indissociables dans un contexte professionnel
où les alertes doivent non seulement être justes, mais aussi justifiées. Ce chapitre
synthétise les résultats issus des notebooks \texttt{03\_baseline\_models.ipynb},
\texttt{04\_autoencoder.ipynb}, \texttt{05\_ollama\_integration.ipynb} et
\texttt{06\_shap\_lime.ipynb}.

% ---------------------------------------------------
\section{Métriques d'évaluation}
% ---------------------------------------------------

\subsection{Choix et justification des métriques}

Dans un contexte de fort déséquilibre, l'exactitude (\textit{accuracy}) n'est pas
pertinente : un modèle qui prédit systématiquement « légitime » obtiendrait 99,87\,\%
d'exactitude tout en ne détectant aucune fraude. Les métriques retenues sont :

\begin{itemize}
  \item \textbf{Rappel (Recall)} : fraction des fraudes réelles correctement détectées.
    \textit{Priorité~1} car manquer une fraude a un coût financier et réputationnel direct.
    \[
    \text{Recall} = \frac{TP}{TP + FN}
    \]
  \item \textbf{Précision (Precision)} : fraction des alertes qui sont réellement des
    fraudes. \textit{Priorité~2} car trop de faux positifs surchargent les équipes
    d'investigation.
    \[
    \text{Precision} = \frac{TP}{TP + FP}
    \]
  \item \textbf{F1-Score} : moyenne harmonique de la précision et du rappel. Synthèse
    unique du compromis.
    \[
    \text{F1} = \frac{2 \cdot \text{Precision} \cdot \text{Recall}}
    {\text{Precision} + \text{Recall}}
    \]
  \item \textbf{PR-AUC} (\textit{Area Under the Precision-Recall Curve}) : synthétise
    le tradeoff précision/rappel sur l'ensemble des seuils possibles --- robuste au
    déséquilibre des classes.
\end{itemize}

\subsection{Résultats comparatifs sur l'ensemble de test}

\begin{table}[H]
\centering
\caption{Comparaison des performances sur l'ensemble de test}
\label{tab:results_comparison}
\begin{tabular}{lccccc}
\toprule
\textbf{Modèle} & \textbf{Seuil} & \textbf{Recall}
  & \textbf{Precision} & \textbf{F1-Score} & \textbf{PR-AUC} \\
\midrule
Baseline (\texttt{isFlaggedFraud}) & N/A    & 0,0039 & 1,0000 & 0,0077 & N/A \\
LR\_balanced                       & 0,9986 & 0,6410 & 0,6579 & 0,6494 & 0,7044 \\
LR\_smote                          & 0,9621 & 0,6923 & 0,7297 & 0,7105 & 0,7393 \\
RF\_balanced                       & 0,6235 & 0,7692 & 0,8333 & 0,8000 & 0,7794 \\
RF\_smote                          & 0,6291 & 0,7949 & 0,8158 & 0,8052 & 0,8405 \\
Autoencoder                        & 1,6874 & 0,3590 & 0,6087 & 0,4516 & 0,3734 \\
\bottomrule
\end{tabular}
\end{table}

\textit{Note : les seuils sont sélectionnés sur la validation, puis figés pour l'évaluation
finale sur l'ensemble de test --- jamais consulté lors de la phase d'optimisation.}

\subsection{Analyse de la matrice de confusion}

La matrice de confusion permet d'analyser précisément les types d'erreurs :
\begin{itemize}
  \item \textbf{Faux négatifs (FN)} : fraudes manquées --- coût financier direct,
    impact réputationnel, responsabilité de l'auditeur.
  \item \textbf{Faux positifs (FP)} : alertes injustifiées --- surcharge des équipes
    d'investigation, coût opérationnel.
\end{itemize}

La forêt aléatoire avec SMOTE offre le meilleur équilibre sur l'ensemble de test :
\textbf{31 fraudes détectées sur 39} pour seulement \textbf{7 faux positifs}. La variante
\texttt{RF\_balanced} reste très proche (30 TP, 6 FP), tandis que \texttt{LR\_smote}
conserve un compromis acceptable (27 TP, 10 FP). L'autoencoder, en revanche, ne détecte
que \textbf{14 fraudes sur 39} avec \textbf{9 faux positifs} --- moins performant que les
approches supervisées, mais conservant un intérêt lorsque les labels sont rares ou
indisponibles.

\subsection{Discussion et sélection du modèle}

Tous les critères de succès définis en phase~1 sont atteints pour la forêt aléatoire avec
SMOTE :

\begin{table}[H]
\centering
\caption{Vérification des critères de succès pour RF\_smote}
\label{tab:criteres_verification}
\begin{tabular}{llll}
\toprule
\textbf{Critère} & \textbf{Cible} & \textbf{Résultat} & \textbf{Statut} \\
\midrule
Recall   & $> 0{,}70$ & 0,7949 & \checkmark Atteint \\
F1-Score & $> 0{,}70$ & 0,8052 & \checkmark Atteint \\
PR-AUC   & $> 0{,}75$ & 0,8405 & \checkmark Atteint \\
\bottomrule
\end{tabular}
\end{table}

% ---------------------------------------------------
\section{Interprétabilité des modèles}
% ---------------------------------------------------

\subsection{Méthodes SHAP}

L'un des reproches fréquents adressés aux modèles d'apprentissage automatique est leur
opacité. Dans un contexte d'audit, un analyste ne peut se contenter d'une alerte : il lui
faut une justification claire des variables qui ont conduit le modèle à signaler une
transaction.

\begin{definitionbox}
\textbf{SHAP} (\textit{SHapley Additive exPlanations}) est une méthode d'interprétabilité
fondée sur la théorie des jeux coopératifs. Elle attribue à chaque variable une valeur
représentant sa contribution marginale à la prédiction, calculée en moyennant l'impact de
la variable sur toutes les combinaisons possibles de features. Cette approche garantit
additivité, cohérence et absence de biais entre les variables.
\end{definitionbox}

Pour la forêt aléatoire, un explainer \textbf{TreeSHAP} est utilisé --- plus rapide que
l'approche agnostique et exact pour les modèles arborescents. Les résultats montrent que
\texttt{balance\_diff\_orig}, \texttt{is\_transfer\_or\_cashout} et
\texttt{log\_amount} sont les variables à plus forte contribution, en cohérence avec
l'analyse exploratoire et l'importance des variables de la forêt.

\textbf{Deux niveaux d'analyse SHAP sont produits :}
\begin{itemize}
  \item \textbf{Global} : beeswarm plot montrant la distribution des valeurs SHAP sur
    l'ensemble du test, permettant d'identifier les variables globalement les plus
    influentes.
  \item \textbf{Local} : waterfall plot pour une transaction spécifique, montrant
    comment chaque feature pousse la prédiction vers « fraude » ou « légitime ».
\end{itemize}

\subsection{Méthodes LIME}

\begin{definitionbox}
\textbf{LIME} (\textit{Local Interpretable Model-agnostic Explanations}) explique chaque
prédiction individuellement en perturbant légèrement la transaction d'entrée et en
observant comment la prédiction varie. Il approxime localement le comportement du modèle
complexe par un modèle linéaire simple, interprétable par construction, valable uniquement
dans le voisinage de l'exemple analysé.
\end{definitionbox}

LIME est particulièrement utile pour expliquer un cas suspect spécifique à un analyste,
sans avoir besoin de comprendre le comportement global du modèle. Contrairement à SHAP,
il est \textbf{agnostique} à l'architecture du modèle et peut s'appliquer à n'importe
quel algorithme --- y compris l'autoencoder, pour lequel SHAP n'est pas nativement adapté.

\subsection{Comparaison SHAP vs LIME}

\begin{table}[H]
\centering
\caption{Comparaison des méthodes d'explicabilité}
\label{tab:shap_lime}
\begin{tabular}{lll}
\toprule
\textbf{Critère} & \textbf{SHAP} & \textbf{LIME} \\
\midrule
Scope          & Global + Local & Local uniquement \\
Fondement      & Théorie des jeux & Approximation locale linéaire \\
Fidélité       & Élevée (TreeSHAP exact) & Variable (approximation) \\
Rapidité       & Rapide (TreeSHAP) & Lent (nombreuses perturbations) \\
Agnosticisme   & Méthodes spécifiques disponibles & Entièrement agnostique \\
\bottomrule
\end{tabular}
\end{table}

% ---------------------------------------------------
\section{Intégration Ollama}
% ---------------------------------------------------

\subsection{Pourquoi un LLM local ?}

Ollama est un outil permettant d'exécuter des grands modèles de langage (LLM)
\textbf{entièrement en local}, sans envoyer de données vers des services cloud externes.
Ce point était non négociable dans le cadre de ce projet : les transactions financières
auditées sont confidentielles et ne peuvent pas quitter l'environnement sécurisé du client.
L'intégration avec Ollama permet de produire des explications en langage naturel ---
directement lisibles par un auditeur sans formation en machine learning --- sans compromettre
la confidentialité des données.

\subsection{Architecture de l'intégration}

Le module \texttt{src/llm\_integration/llm\_helper.py} encapsule la logique d'appel au
modèle de langage. Pour chaque transaction signalée comme anormale, les éléments suivants
sont transmis au LLM dans un prompt structuré :

\begin{itemize}
  \item Le score d'anomalie ou d'erreur de reconstruction de l'autoencoder ;
  \item Les valeurs des variables les plus contributives (issues de SHAP/LIME) ;
  \item Le contexte métier de la transaction (type, montant, heure, différence de solde).
\end{itemize}

Le LLM génère une explication structurée en quatre points : type de transaction à risque,
heure suspecte, comportement anormal des soldes, et recommandation d'action.

\subsection{Difficultés rencontrées}

\begin{itemize}
  \item \textbf{Instabilité des seuils de l'autoencoder} : de légères variations du seuil
    peuvent significativement modifier le ratio précision/rappel, nécessitant une
    optimisation soigneuse sur la validation.
  \item \textbf{Variabilité des explications LLM} : les explications Ollama peuvent varier
    selon la version du modèle utilisé et la qualité du prompt ; un prompt engineering
    rigoureux a été nécessaire.
  \item \textbf{Temps de calcul SHAP} : le calcul des valeurs SHAP pour l'ensemble de
    test complet nécessite plusieurs minutes sur CPU, ce qui a motivé la sauvegarde des
    tenseurs SHAP dans \texttt{outputs/models/shap\_values\_*.npy}.
\end{itemize}

% ---------------------------------------------------
\section*{Conclusion}
\addcontentsline{toc}{section}{Conclusion}
% ---------------------------------------------------

L'évaluation confirme que la forêt aléatoire avec SMOTE constitue la \textbf{meilleure
approche supervisée} pour ce jeu de données, avec un recall de 0,7949, un F1-score de
0,8052 et un PR-AUC de 0,8405 --- tous au-dessus des cibles fixées. L'autoencoder apporte
une perspective complémentaire non supervisée, utile lorsque les étiquettes sont limitées.
La \textbf{double couche d'explicabilité} (SHAP/LIME + Ollama) répond à l'objectif
d'interprétabilité en rendant les prédictions compréhensibles tant pour les data scientists
que pour les analystes métier non techniciens. Le chapitre suivant présente les aspects
de déploiement et les garanties de reproductibilité du pipeline complet.
